% ==================================================================
% id:       prob-007
% topic:    מבוא להסתברות
% subtopic: מרחב הסתברות אחיד
% solution: full
% ==================================================================
מטילים שתי קוביות הוגנות, אחת אחרי השנייה. מרחב המדגם הוא $\Omega=\set{1,\dots,6}\times\set{1,\dots,6}$, ו-$|\Omega|=36$. בכל סעיף, כתבו את המאורע כתת-קבוצה של $\Omega$, ומצאו את ההסתברות שלו.
\begin{enumerate}
\item סכום התוצאות הוא 4.
\begin{solbox}
$A=\set{(1,3),(2,2),(3,1)}$, ולכן $\PP(A)=\frac3{36}=\ans{\frac1{12}}$.
\end{solbox}
\item התוצאה בקובייה הראשונה גדולה לפחות ב-4 מהתוצאה בקובייה השנייה.
\begin{solbox}
צריך $a\ge b+4$. אם $b=1$ אז $a\in\set{5,6}$, ואם $b=2$ אז $a=6$. עבור $b\ge3$ אין אפשרויות. לכן $B=\set{(5,1),(6,1),(6,2)}$, ו-$\PP(B)=\frac3{36}=\ans{\frac1{12}}$.
\end{solbox}
\item סכום התוצאות קטן מ-13.
\begin{solbox}
הסכום תמיד בין 2 ל-12, ולכן המאורע הוא כל $\Omega$, ו-$\PP(\Omega)=\ans1$.
\end{solbox}
\item סכום התוצאות הוא 1.
\begin{solbox}
הסכום תמיד לפחות 2, ולכן המאורע הוא $\emptyset$, ו-$\PP(\emptyset)=\ans0$.
\end{solbox}
\end{enumerate}
